\section{System Overview}
\label{sec:overview}

\subsection{Problem Formulation}
\label{sec:problem}

\textbf{Estimation problem.}
We estimate the continuous-time 6-DoF body trajectory (world-frame
position $\vp_w(t)\in\mathbb{R}^3$ and orientation $\vR(t)\in\SO$)
together with the constant \ac{imu} biases $\vb_a,\vb_g$, from asynchronous
measurements over the estimation interval (the full trajectory in batch
mode, a sliding window online): radar Doppler velocities ($\sim$11\,Hz),
accelerometer and gyroscope samples (1\,kHz), and a low-rate heading
prior.  We adopt a \emph{continuous-time} representation because the
sensors are asynchronous and run at disparate rates: a trajectory defined
for every $t$ is evaluated at each measurement's exact timestamp, and its
derivatives (velocity $\dot{\vp}_w$, acceleration $\ddot{\vp}_w$, and
body angular velocity $\vom_b$) follow in closed form, so no
inter-sample interpolation or rigid-keyframe approximation is required.

\textbf{Parameterization.}
The trajectory is a B-spline (\autoref{sec:state}): the position is a
quintic uniform B-spline with control points
$\{\mathbf{P}_i\}\subset\mathbb{R}^3$, and the orientation a cumulative
$\SO$ B-spline with \emph{absolute} rotation knots
$\{\vR_i\}\subset\SO$ (basalt convention~\cite{usenko2020basalt}).
This renders the otherwise infinite-dimensional trajectory
finite-dimensional; the control points and the \ac{imu} biases form the
estimated state
\begin{equation}
  \mathcal{X} = \bigl\{\,\mathbf{P}_i,\ \vR_i,\ \vb_a,\ \vb_g\,\bigr\}.
  \label{eq:state}
\end{equation}
Only the radar pitch is observable among the extrinsics; the
sliding window fixes it at the measured value (self-calibration analyzed in
\autoref{sec:backflips}).

\textbf{Estimator.}
Treating the measurement noise as zero-mean and independent, the
\ac{map} estimate is the minimizer of the negative
log-posterior, a weighted nonlinear least-squares problem over the
measurement residuals, the spline regularizers, and the priors:
\begin{equation}
\begin{aligned}
  \hat{\mathcal{X}} = \operatorname*{arg\,min}_{\mathcal{X}}
    & \sum_{k}\rho_\delta\!\bigl(r_{D,k}^2\bigr)
    + \lambda_a\!\sum_{m}\lVert r_{a,m}\rVert^2
    + \lambda_g\!\sum_{m}\lVert r_{g,m}\rVert^2 \\
    & + \lambda_\psi\!\sum_{n} r_{\psi,n}^2
    + E_{\text{reg}}(\mathcal{X})
    + E_{\text{prior}}(\mathcal{X}),
\end{aligned}
\label{eq:map}
\end{equation}
where $r_D,r_a,r_g,r_\psi$ are the radar-Doppler, accelerometer,
gyroscope, and heading residuals of \autoref{sec:models};
$\rho_\delta$ is the Huber kernel ($\delta{=}1.0$\,m/s) that bounds the
influence of heavy-tailed Doppler outliers; the $\lambda_\bullet$ are the
(rate-adaptive) information weights of \autoref{sec:models};
$E_{\text{reg}}$ collects the position min-snap and orientation
smoothness regularizers (\autoref{sec:reg}); and $E_{\text{prior}}$
holds the boundary and, online, the marginalization priors
(\autoref{sec:sw}).  We minimize~\eqref{eq:map} with
Levenberg--Marquardt (Ceres) using analytic Jacobians.  The fixed-lag
estimator solves it over a 3\,s window and, on each stride, marginalizes
the states leaving the window into $E_{\text{prior}}$ via the Schur
complement (\autoref{sec:sw}), keeping each window's problem size bounded.

\subsection{Hardware Platform}
The system runs on a quadrotor with a TI IWR6843AOPEVM
mmWave radar (60--64\,GHz \ac{fmcw}), a Pixhawk flight
controller \ac{imu}, and a Vicon \ac{mocap} system used for ground
truth evaluation.  The radar is mounted upside-down; the extrinsic
rotation is $[180\degree, 27.5\degree, 0\degree]$
(roll, pitch, yaw) and translation $[0.08, 0.02, -0.01]$\,m in the body
frame, with the $27.5\degree$ pitch the as-built downtilt measured by
inclinometer (\autoref{sec:results}).  The radar boresight lies along the body forward axis, so the
Euler pitch is numerically the physical boresight downtilt.  The radar operates at $\sim$11\,Hz with 0.049\,m/s Doppler
resolution (best-velocity firmware configuration).  Raw \ac{imu} rate
is $\sim$993\,Hz; both modalities are used at full rate.

\subsection{Pipeline}
\autoref{fig:pipeline} shows the pipeline.  A three-stage sensor-only
initialization (P1--P3, detailed in \autoref{sec:init}) builds the
initial trajectory the nonlinear solve needs, from radar and \ac{imu}
data alone; the same raw measurements then also enter the solver directly
as factors (\autoref{fig:pipeline}, dashed).  \ac{mocap} supplies only the
yaw heading prior (a noise-free, idealized magnetometer surrogate; \autoref{sec:yaw} bounds the real-sensor gap) and the
final \ac{rmse} evaluation; P1--P3 and the position/attitude estimate use
no external pose reference.

\begin{figure}[t]
\centering
\definecolor{loopblue}{RGB}{31,119,180}
\resizebox{\columnwidth}{!}{%
\begin{tikzpicture}[
  font=\normalsize,
  box/.style={rectangle, rounded corners=3pt, draw, minimum width=1.55cm,
              minimum height=0.72cm, align=center, inner sep=2pt},
  sensor/.style={box, fill=blue!12},
  stage/.style={box, fill=orange!22},
  solver/.style={box, fill=green!16, minimum width=1.7cm, minimum height=1.05cm},
  outb/.style={box, fill=gray!12, minimum width=1.25cm},
  arr/.style={-{Latex[length=2.2mm]}, thick},
  darr/.style={-{Latex[length=2.2mm]}, thick, dashed, black!55},
  loop/.style={-{Latex[length=2.0mm]}, thick, loopblue},
  lbl/.style={font=\small},
  ]
  \node[sensor] (imu)   at (0.0,  1.6)  {\ac{imu}\\1\,kHz};
  \node[sensor] (radar) at (0.0,  0.35) {Radar\\$\sim$11\,Hz};
  \node[stage]  (p1)    at (2.35,  1.6) {P1: Gyro\\${\to}\,R(t)$};
  \node[stage]  (p2)    at (2.35,  0.35) {P2: \ac{wls}\\${\to}\,\mathbf{p}(t)$};
  \node[stage]  (p3)    at (4.35,  1.0) {P3: Priors\\${\to}$ init};
  \node[solver] (solver) at (6.35, 1.0) {Ceres\\LM solve\\{\small(batch / SW)}};
  \node[outb]   (out)   at (8.2,  1.0) {$\hat{\mathbf{p}},\,\hat{R},\,\hat{\mathbf{v}}$};
  \begin{scope}[on background layer]
    \node[draw=black!40, dashed, rounded corners=5pt, fill=orange!6,
          inner sep=6pt, fit=(p1)(p2)(p3)] (initbox) {};
    \node[lbl, black!65, anchor=south] at (initbox.north)
          {Initialization (once)};
  \end{scope}
  \draw[arr] (imu.east)   -- (p1.west);
  \draw[arr] (imu.east)   -- (p2.west);
  \draw[arr] (radar.east) -- (p2.west);
  \draw[arr] (p1.east)    -- (p3.north west);
  \draw[arr] (p2.east)    -- (p3.south west);
  \draw[arr] (p3.east)    -- (solver.west);
  \draw[arr] (solver.east)-- (out.west);
  \coordinate (busY) at (0,-0.95);
  \draw[darr] (imu.west) -- ++(-0.32,0) |- (busY) -| (solver.south);
  \draw[darr] (radar.west) -- ++(-0.5,0) |- ($(busY)+(0,-0.42)$)
        -| ($(solver.south)+(0.25,0)$);
  \draw[loop] ([xshift=-3mm]solver.north) to[out=115,in=65,looseness=3.2]
        node[lbl, loopblue, above=0pt] {per stride}
        ([xshift=3mm]solver.north);
\end{tikzpicture}%
}% end resizebox
\caption{System pipeline. Solid: initialization flow (P1--P3 seed
the optimizer, run once). Dashed: raw sensor data fed as measurement
factors. Blue: the sliding-window re-solve at each stride.
Radar returns first pass a \ac{ransac} ego-velocity prefilter
(\autoref{sec:models}) that rejects elevation-biased outliers before
both P2 and the solver. \ac{mocap} supplies only the yaw heading prior
(magnetometer surrogate) and \ac{rmse} evaluation.}
\label{fig:pipeline}
\end{figure}
